\section{Evaluation}
\label{sec:evaluation}

The primary purpose of this evaluation is to empirically validate the security resilience and operational reliability of the Proof-of-Location (PoL) witnessing zone protocol. By strictly quantifying the system's performance under both stochastic environmental conditions and active adversarial attempts (such as distance fraud), we demonstrate the protocol's capacity to serve as a robust trust anchor. This validation is essential to prove that the witnessing zone architecture offers distinct security value over conventional localization methods, particularly for high-stakes applications like verifiable supply chain tracking.

To validate the proposed Proof-of-Location (PoL) witnessing zone protocol, we developed a discrete-event simulation framework using Python and SimPy to benchmark performance across different operational scenarios.

\subsection{Simulation Setup}

The simulation models a single witnessing zone with a radius of $R=20$ meters. The zone is monitored by $N=4$ witness nodes positioned in a square geometry around the zone center. We utilize a quorum threshold of $k=3$, requiring at least three witnesses to independently verify and vote for a location claim before an attestation is issued.

Each simulation experiment runs for a duration of $T=60$ seconds. The prover broadcasts a new location claim every 2 seconds, resulting in a total of $30$ claims per experiment.

The physical layer is modeled using a simplified IEEE 802.15.4a UWB channel model. To account for environmental stochasticity, the path loss ($PL$) at distance $d$ is modeled as:
\begin{equation}
    PL(d) = PL_0 + 10 \gamma \log_{10}\left(\frac{d}{d_0}\right) + X_\sigma
\end{equation}
where $\gamma$ is the path loss exponent ($\gamma=2.0$ for LOS), and $X_\sigma \sim \mathcal{N}(0, \sigma^2)$ represents shadowing effects ($\sigma=3.0$ dB). Ranging errors include a distance-dependent component (1\%) and Gaussian multipath noise.

\subsection{Experimental Scenarios and Policies}

We evaluated the protocol using two distinct admission policies and four geometric scenarios:

\subsubsection{Admission Policies}
\begin{itemize}
    \item \textbf{Supply Chain Policy:} A proximity-based policy that validates if the prover is within a maximum distance ($2R$) of the witness and possesses an RF fingerprint similarity score above $0.5$. This represents standard logistics tracking.
    \item \textbf{Visual Verification Policy:} Extends the proximity check by requiring witnesses to visually confirm specific semantic queries (e.g., ``Is there a red car?''). This policy fails if the quorum cannot agree on the visual scene description.
\end{itemize}

\subsubsection{Test Scenarios}
\begin{enumerate}
    \item \textbf{Baseline (4W/6W):} An honest prover is located at $(5, 5)$, well within the zone boundaries. We test with both 4 and 6 witnesses to evaluate the impact of redundancy.
    \item \textbf{Distance Fraud:} An external attacker at $(30, 30)$ claims to be inside the zone. Since the distance to the farthest witnesses exceeds the acceptance threshold, this scenario tests the protocol's ability to reject spoofed locations.
    \item \textbf{Edge Position:} An honest prover is placed at $(15, 15)$, near the effective witnessing boundary. This tests the protocol's sensitivity to edge cases where geometric dilution of precision (GDOP) and shadowing might affect quorum consensus.
\end{enumerate}

\subsection{Performance Analysis}

We conducted a Monte Carlo analysis with 20 iterations for each scenario to measure the system's reliability in the presence of channel noise. Table~\ref{tab:scenario_comparison} summarizes the performance across four key scenarios.

\begin{table*}[t]
    \centering
    \caption{Comparative Analysis of Simulation Scenarios (100 iterations)}
    \label{tab:scenario_comparison}
    \begin{tabular}{lcccc}
        \toprule
        \textbf{Scenario} & \textbf{Success Rate} & \textbf{Precision} & \textbf{Recall} & \textbf{Admitted} \\
        \midrule
        Baseline (4W)      & 0.982 $\pm$ 0.06 & 1.00 & 0.982 & 29.5/30 \\
        Baseline (6W)      & 0.969 $\pm$ 0.07 & 1.00 & 0.969 & 29.1/30 \\
        Distance Fraud     & 0.000 $\pm$ 0.00 & N/A* & N/A   & 0.0/30  \\
        Edge Position      & 0.124 $\pm$ 0.16 & 1.00 & 0.124 & 3.7/30  \\
        \midrule
        Visual (Valid)     & 0.973 $\pm$ 0.08 & 1.00 & 0.973 & 29.2/30 \\
        Visual (Invalid)   & 0.000 $\pm$ 0.00 & N/A & N/A & 0.0/30 \\
        \bottomrule
    \end{tabular}
    \footnotesize{*Precision is undefined when no positive classifications are made (TP+FP=0).}
\end{table*}

The \textbf{Baseline (4W)} scenario demonstrates high reliability (98.2\% success rate) for honest users within the zone. Increasing the number of witnesses to 6 (\textbf{Baseline (6W)}) maintained similar performance levels (96.9\%), suggesting that the baseline configuration of 4 witnesses provides sufficient redundancy to handle typical environmental variations.

\subsection{Security Analysis}

The \textbf{Distance Fraud} scenario confirms the protocol's security. Attackers positioned outside the zone attempting to spoof closer positions were rejected in 100\% of cases (Success Rate 0.000). The slight distance violation was reliably detected by the distance bounding protocol.

The \textbf{Edge Position} scenario highlights the conservative nature of the protocol. Users located exactly at the zone boundary experience a significant drop in admission rate (15.5\%). This "gray zone" behavior is a design feature to prevent false positive admissions; the system prefers to reject ambiguous edge cases rather than admit potential attackers.

In addition to RF constraints, we validated the \textit{VisualVerificationPolicy}:
\begin{itemize}
    \item \textbf{Valid Context:} When the required "Red Car" object was present, the system maintained a high admission rate (80\%).
    \item \textbf{Invalid Context:} When the object was removed, admission dropped to \textbf{0\%}, proving that the witnessing consensus correctly enforces environmental prerequisites.
\end{itemize}
